% Root candidate. No existence of a suitable odd-degree neighbour is asserted.
\subsection{Triple fibres are the only obstruction for a nonordinary theta norm}

\begin{lemma}[Odd theta norms in two-rank zero or one]
\label{lem:root-nonordinary-odd-theta-nodality}
\status{Candidate; submitted for independent proof audit.}
Let $C$ be a smooth projective genus-two curve over an algebraically
closed field $k$ of characteristic two, with two-rank $f\leq1$.
Choose a Weierstrass point $W$ and identify
$C=\Theta\subset B=\operatorname{Jac}(C)$ by $P\mapsto[P-W]$.
Let $K\subset B(k)$ be a finite subgroup of odd order, and let
$\phi:B\to A=B/K$ be the quotient isogeny.
Then $\phi|_C$ is birational onto its integral image and is an
immersion at each point. Its image is nodal if and only if every
geometric fibre of $\phi|_C$ has at most two points. In that case
the image has exactly $|K|-1$ nodes.
\end{lemma}
\begin{proof}
The isogeny is etale, so the restriction of its differential to
each tangent line of $C$ is injective. A nonzero translation cannot
stabilize $\Theta$: it would belong to the kernel of the principal
polarization. For each nonzero $h\in K$, the distinct curves
$\Theta$ and $\Theta+h$ meet in finitely many points. Away from
the union of these finitely many intersections, the restriction
of $\phi$ is one-to-one. Thus it is birational, and the smooth
curve $C$ is the normalization of its integral image.

We prove that any nontransverse pair of distinct branches forces
a third point in the same fibre. Restriction of invariant
differentials identifies
$H^0(B,\Omega_B^1)$ with $H^0(C,\omega_C)$. Hence equality of
two translated tangent lines is equality of their images under
the canonical map of $C$. Its degree-two fibres are hyperelliptic
pairs. A nontransverse pair therefore comes from $P$ and
$\iota P$, with $P\ne\iota P$.
Under the chosen Abel--Jacobi embedding these have classes $x$
and $-x$, so their common image $y=\phi(x)$ satisfies $2y=0$.

The map $\phi:B[2](k)\to A[2](k)$ is bijective. For example,
if $2\phi(x)=0$, then $2x\in K$; multiplication by two is
bijective on $K$, so subtracting its unique half gives a
two-torsion preimage of $\phi(x)$. Its uniqueness follows from
$K\cap B[2](k)=0$.
For $f=0$, $B[2](k)=\{0\}$, represented on $C$ by $W$.
For $f=1$, the two Weierstrass points $W,W'$ give exactly
$B[2](k)=\{0,[W'-W]\}$, as in the previously proved
two-rank-one normal form; here $2W\sim2W'\sim K_C$.
Consequently the two-torsion preimage of $y$ is represented
by a Weierstrass point $Q$ of $C$. It is distinct from both
$P$ and $\iota P$, because those are not fixed by $\iota$.
The fibre over $y$ thus has at least three points.

It follows that a fibre with two points has transverse tangent
lines. The etale local nature of $\phi$ makes a one-point fibre
smooth and a transverse two-point fibre an ordinary node.
Conversely, a node has exactly two points in its normalization,
so a fibre with at least three points is not nodal.

Finally, putting $D=\phi(C)$, we have the divisor identity
$\phi^*D=\sum_{h\in K}(\Theta+h)$, with distinct components
and multiplicity one. Projection formula gives
$|K|D^2=2|K|^2$, hence $D^2=2|K|$.
Adjunction on the abelian surface gives $p_a(D)=|K|+1$.
Its normalization has genus two, so when all singularities
are nodes their number is $p_a(D)-2=|K|-1$.
\end{proof}

\begin{remark}[The incoming isogeny remains an existence requirement]
\label{rem:root-nonordinary-odd-theta-nodality-scope}
\status{Candidate scope statement; no new degree is classified here.}
For the proved odd norm-transfer theorem it now suffices to find
an incoming polarized odd isogeny whose theta curve has no
three-point fibre. This lemma supplies the tangency exclusion;
it does not supply that isogeny. In particular, finding a good
outgoing quotient of a fixed surface does not by itself give
the required incoming neighbour of that same surface.
Neither a degree-five neighbour for every fixed two-rank-one
surface nor a triple-free degree-seven self-isogeny is asserted.
\end{remark}
